\documentclass[11pt]{article}
% The preamble every project's paper shares: % The preamble every project's paper shares: % The preamble every project's paper shares: \input{../../preamble} from
% projects/<project>/paper/. Boards come from chessboard and xskak, plots from
% pgfplots; every number a paper prints comes from its generated.tex.
\usepackage[LSBC4,T1]{fontenc}
\usepackage[utf8]{inputenc}
\usepackage{lmodern}
\usepackage[margin=1in]{geometry}
\usepackage{microtype}
\usepackage{amsmath}
\usepackage{booktabs}
\usepackage{tabularx}
\usepackage{array}
\usepackage{xcolor}
\usepackage[obeyspaces,spaces]{url}
\usepackage{xskak}
\usepackage{pgfplots}
\pgfplotsset{compat=1.18}
\usepackage{caption}
\usepackage{subcaption}
\usepackage{enumitem}
\usepackage[section]{placeins}
\usepackage[hidelinks]{hyperref}

\setchessboard{boardfontencoding=LSBC4, boardfontfamily=skaknew, blackfieldcolor=black!22,
  whitefieldcolor=white, setfontcolors, showmover=false, label=true, labelfontsize=6pt,
  boardfontsize=14pt}

% A move or line in UCI, and a position in FEN, both breakable at a slash.
\newcommand{\uci}[1]{\texttt{#1}}
\newcommand{\fen}[1]{{\footnotesize\path{#1}}}
\newcolumntype{L}{>{\raggedright\arraybackslash}X}

\setlength{\parskip}{0.4em}
\setlength{\parindent}{0pt}
 from
% projects/<project>/paper/. Boards come from chessboard and xskak, plots from
% pgfplots; every number a paper prints comes from its generated.tex.
\usepackage[LSBC4,T1]{fontenc}
\usepackage[utf8]{inputenc}
\usepackage{lmodern}
\usepackage[margin=1in]{geometry}
\usepackage{microtype}
\usepackage{amsmath}
\usepackage{booktabs}
\usepackage{tabularx}
\usepackage{array}
\usepackage{xcolor}
\usepackage[obeyspaces,spaces]{url}
\usepackage{xskak}
\usepackage{pgfplots}
\pgfplotsset{compat=1.18}
\usepackage{caption}
\usepackage{subcaption}
\usepackage{enumitem}
\usepackage[section]{placeins}
\usepackage[hidelinks]{hyperref}

\setchessboard{boardfontencoding=LSBC4, boardfontfamily=skaknew, blackfieldcolor=black!22,
  whitefieldcolor=white, setfontcolors, showmover=false, label=true, labelfontsize=6pt,
  boardfontsize=14pt}

% A move or line in UCI, and a position in FEN, both breakable at a slash.
\newcommand{\uci}[1]{\texttt{#1}}
\newcommand{\fen}[1]{{\footnotesize\path{#1}}}
\newcolumntype{L}{>{\raggedright\arraybackslash}X}

\setlength{\parskip}{0.4em}
\setlength{\parindent}{0pt}
 from
% projects/<project>/paper/. Boards come from chessboard and xskak, plots from
% pgfplots; every number a paper prints comes from its generated.tex.
\usepackage[LSBC4,T1]{fontenc}
\usepackage[utf8]{inputenc}
\usepackage{lmodern}
\usepackage[margin=1in]{geometry}
\usepackage{microtype}
\usepackage{amsmath}
\usepackage{booktabs}
\usepackage{tabularx}
\usepackage{array}
\usepackage{xcolor}
\usepackage[obeyspaces,spaces]{url}
\usepackage{xskak}
\usepackage{pgfplots}
\pgfplotsset{compat=1.18}
\usepackage{caption}
\usepackage{subcaption}
\usepackage{enumitem}
\usepackage[section]{placeins}
\usepackage[hidelinks]{hyperref}

\setchessboard{boardfontencoding=LSBC4, boardfontfamily=skaknew, blackfieldcolor=black!22,
  whitefieldcolor=white, setfontcolors, showmover=false, label=true, labelfontsize=6pt,
  boardfontsize=14pt}

% A move or line in UCI, and a position in FEN, both breakable at a slash.
\newcommand{\uci}[1]{\texttt{#1}}
\newcommand{\fen}[1]{{\footnotesize\path{#1}}}
\newcolumntype{L}{>{\raggedright\arraybackslash}X}

\setlength{\parskip}{0.4em}
\setlength{\parindent}{0pt}

% Written by projects/lesson-moves/paper/generate.mjs; do not edit by hand.

\newcommand{\AppCommit}{f7a8b12a01}
\newcommand{\Lessons}{74}
\newcommand{\Chapters}{527}
\newcommand{\TrainUntil}{2026-09-16}
\newcommand{\HeldoutUntil}{2026-11-02}
\newcommand{\Draws}{10,000}
\newcommand{\CoverageBar}{50}
\newcommand{\CostBar}{1.5}
\newcommand{\Window}{3}
\newcommand{\KeptWon}{88}
\newcommand{\Offset}{250}
\newcommand{\Examined}{1,511}
\newcommand{\SkippedGames}{seven}
\newcommand{\MaxCandidates}{five}
\newcommand{\TrainTPs}{1,023}
\newcommand{\TrainGames}{399}
\newcommand{\HeldTPs}{488}
\newcommand{\HeldGames}{191}
\newcommand{\TrainCoverage}{77.3}
\newcommand{\TrainNoLesson}{22.7}
\newcommand{\TrainFresh}{63.7}
\newcommand{\TrainCost}{0.32 (0.26 to 0.38)}
\newcommand{\TrainCostMean}{0.32}
\newcommand{\TrainWorst}{10.2}
\newcommand{\TrainFind}{3.0 (2.2 to 3.9)}
\newcommand{\TrainFindMean}{3.0}
\newcommand{\TrainVsNatural}{$-$2.8 ($-$3.5 to $-$2.2)}
\newcommand{\TrainIsTop}{74}
\newcommand{\TrainIsNatural}{87}
\newcommand{\TrainTopTaught}{68}
\newcommand{\TrainLessonsNamed}{20}
\newcommand{\TrainChaptersNamed}{25}
\newcommand{\HeldCoverage}{71.1}
\newcommand{\HeldNoLesson}{28.9}
\newcommand{\HeldFresh}{58.0}
\newcommand{\HeldCost}{0.37 (0.29 to 0.46)}
\newcommand{\HeldCostMean}{0.37}
\newcommand{\HeldWorst}{11.1}
\newcommand{\HeldFind}{3.6 (2.5 to 4.8)}
\newcommand{\HeldFindMean}{3.6}
\newcommand{\HeldVsNatural}{$-$2.3 ($-$3.0 to $-$1.6)}
\newcommand{\HeldIsTop}{74}
\newcommand{\HeldIsNatural}{87}
\newcommand{\HeldTopTaught}{59}
\newcommand{\HeldLessonsNamed}{21}
\newcommand{\HeldChaptersNamed}{25}
\newcommand{\RuleTable}{%
\begin{table}[ht]\centering\small
\begin{tabular}{lrr}
\toprule
 & Train & Held out \\
\midrule
Turning points (games) & 1,023 (399) & 488 (191) \\
\textbf{Taught by a lesson} (bar: 50\%) & \textbf{77.3\%} & \textbf{71.1\%} \\
\quad for a purpose his move lacked & 63.7\% & 58.0\% \\
\textbf{Loss against E}, win\% (bar: 1.5) & \textbf{0.32 (0.26 to 0.38)} & \textbf{0.37 (0.29 to 0.46)} \\
\quad largest; picks outside the holding verdicts & 10.2; 0 & 11.1; 0 \\
\textbf{Share of L minus share of E}, points (bar: above 0) & \textbf{3.0 (2.2 to 3.9)} & \textbf{3.6 (2.5 to 4.8)} \\
\bottomrule
\end{tabular}
\caption{The rule's three measures on each split. Intervals are 95\% percentiles of 10,000 bootstrap draws over games. A share is Maia-2's probability for the move at his band, so a point is one percent of his band's choices.}
\label{tab:rule}
\end{table}
}
\newcommand{\KeepsPicks}{41}
\newcommand{\OverWindow}{nine}
\newcommand{\TrainDiffers}{261}
\newcommand{\HeldDiffers}{127}
\newcommand{\TrainGainWhere}{11.9 (8.7 to 15.2)}
\newcommand{\HeldGainWhere}{14.0 (10.0 to 18.4)}
\newcommand{\ThreatShare}{36}
\newcommand{\Purposes}{17}
\newcommand{\TaughtN}{1,138}
\newcommand{\ChaptersAll}{26}
\newcommand{\LessonsAll}{21}
\newcommand{\FirstShare}{39}
\newcommand{\TwoShare}{55}
\newcommand{\FirstTitle}{Choose a defense that changes the threat}
\newcommand{\FirstIdea}{Move the queen while taking the loose pawn}
\newcommand{\SecondTitle}{Give your rook an entry square}
\newcommand{\SecondIdea}{Move the rook that has no job}
\newcommand{\FirstNotQueen}{77}
\newcommand{\PurposeTable}{%
\begin{table}[ht]\centering\small
\begin{tabular}{lrr}
\toprule
The pick & Train & Held out \\
\midrule
meets a threat & 285 & 122 \\
a rook to an open file & 111 & 70 \\
wins material & 127 & 46 \\
develops & 59 & 16 \\
mates & 47 & 10 \\
a pawn break & 29 & 19 \\
trades queens & 26 & 18 \\
improves the worst piece & 23 & 8 \\
offers a queen trade & 17 & 10 \\
pushes a passed pawn & 15 & 7 \\
castles & 17 & 5 \\
takes the centre & 12 & 4 \\
an ending technique & 7 & 4 \\
a rook to the seventh rank & 5 & 5 \\
takes an outpost & 6 & 1 \\
trades when ahead & 5 & 1 \\
trades into a pawn ending & 0 & 1 \\
\midrule
Taught by a lesson, all purposes & 791 & 347 \\
\bottomrule
\end{tabular}
\caption{The purpose a lesson teaches the pick for, as the board code reads it from the move and its engine line. These counts describe; they carry no rule.}
\label{tab:purposes}
\end{table}
}
\newcommand{\EaSide}{Black}
\newcommand{\EaFen}{\fen{2r2rk1/pp4pp/5p2/2qp1n2/P7/1P3PP1/3Q2KP/4RR2 b - - 4 31}}
\newcommand{\EaPick}{31\ldots{}\symqueen{}c2}
\newcommand{\EaPickUci}{\uci{c5c2}}
\newcommand{\EaTop}{31\ldots{}d4}
\newcommand{\EaTopUci}{\uci{d5d4}}
\newcommand{\EaPlayed}{31\ldots{}\symqueen{}a3}
\newcommand{\EaPlayedUci}{\uci{c5a3}}
\newcommand{\EaPShare}{40}
\newcommand{\EaTShare}{24}
\newcommand{\EaLoss}{0.8}
\newcommand{\EaPurpose}{trades queens}
\newcommand{\EaSays}{its line trades queens, and with fewer pieces the extra material counts for more}
\newcommand{\EaTitle}{Trade for a reason and inspect what remains}
\newcommand{\EaIdea}{Clear the passer’s path and offer the queens}
\newcommand{\EaBoard}{%
\chessboard[setfen={2r2rk1/pp4pp/5p2/2qp1n2/P7/1P3PP1/3Q2KP/4RR2 b - - 4 31}, boardfontsize=11pt, pgfstyle=straightmove, arrow=to, linewidth=0.07em, color=blue!60!black, markmoves={d5-d4}, pgfstyle=straightmove, arrow=to, linewidth=0.07em, color=green!50!black, markmoves={c5-c2}]
}
\newcommand{\EbSide}{White}
\newcommand{\EbFen}{\fen{3r1rk1/pp3ppp/1n6/3p4/4p2N/1PP5/P4PPP/3RR1K1 w - - 1 22}}
\newcommand{\EbPick}{22.~\symknight{}f5}
\newcommand{\EbPickUci}{\uci{h4f5}}
\newcommand{\EbTop}{22.~g4}
\newcommand{\EbTopUci}{\uci{g2g4}}
\newcommand{\EbPlayed}{22.~c4}
\newcommand{\EbPlayedUci}{\uci{c3c4}}
\newcommand{\EbPShare}{66}
\newcommand{\EbTShare}{under 1}
\newcommand{\EbLoss}{1.1}
\newcommand{\EbPurpose}{improves the worst piece}
\newcommand{\EbSays}{it brings the knight on h4, the piece with the fewest moves (3), to f5, where it has 8}
\newcommand{\EbTitle}{Improve the piece with the weakest job}
\newcommand{\EbIdea}{A central knight needs something to do}
\newcommand{\EbBoard}{%
\chessboard[setfen={3r1rk1/pp3ppp/1n6/3p4/4p2N/1PP5/P4PPP/3RR1K1 w - - 1 22}, boardfontsize=11pt, pgfstyle=straightmove, arrow=to, linewidth=0.07em, color=blue!60!black, markmoves={g2-g4}, pgfstyle=straightmove, arrow=to, linewidth=0.07em, color=green!50!black, markmoves={h4-f5}]
}
\newcommand{\EcSide}{White}
\newcommand{\EcFen}{\fen{r3k2r/pppbqpb1/2n1p2p/3p2p1/3P2P1/2PBPN2/PP1NQPP1/2KR3R w kq - 4 13}}
\newcommand{\EcPick}{13.~e4}
\newcommand{\EcPickUci}{\uci{e3e4}}
\newcommand{\EcTop}{13.~\symknight{}b3}
\newcommand{\EcTopUci}{\uci{d2b3}}
\newcommand{\EcPlayed}{13.~\symknight{}e5}
\newcommand{\EcPlayedUci}{\uci{f3e5}}
\newcommand{\EcPShare}{35}
\newcommand{\EcTShare}{7}
\newcommand{\EcLoss}{0.5}
\newcommand{\EcPurpose}{a pawn break}
\newcommand{\EcSays}{a pawn break: the e-pawn strikes at d5, and trades there open lines}
\newcommand{\EcTitle}{Read the pawn structure before choosing a plan}
\newcommand{\EcIdea}{Attack the supporting pawn}
\newcommand{\EcBoard}{%
\chessboard[setfen={r3k2r/pppbqpb1/2n1p2p/3p2p1/3P2P1/2PBPN2/PP1NQPP1/2KR3R w kq - 4 13}, boardfontsize=11pt, pgfstyle=straightmove, arrow=to, linewidth=0.07em, color=blue!60!black, markmoves={d2-b3}, pgfstyle=straightmove, arrow=to, linewidth=0.07em, color=green!50!black, markmoves={e3-e4}]
}
\newcommand{\EdSide}{White}
\newcommand{\EdFen}{\fen{6k1/1p3p2/p5p1/7p/8/5QPP/Pbq1rP2/3R2K1 w - - 3 30}}
\newcommand{\EdPick}{30.~\symrook{}d8+}
\newcommand{\EdPickUci}{\uci{d1d8}}
\newcommand{\EdTop}{30.~\symrook{}d7}
\newcommand{\EdTopUci}{\uci{d1d7}}
\newcommand{\EdPlayed}{30.~\symking{}f1}
\newcommand{\EdPlayedUci}{\uci{g1f1}}
\newcommand{\EdPShare}{82}
\newcommand{\EdTShare}{4}
\newcommand{\EdLoss}{0.1}
\newcommand{\EdPurpose}{meets a threat}
\newcommand{\EdSays}{it moves the rook on d1 out of attack, to d8}
\newcommand{\EdTitle}{Choose a defense that changes the threat}
\newcommand{\EdIdea}{Move the queen while taking the loose pawn}
\newcommand{\EdBoard}{%
\chessboard[setfen={6k1/1p3p2/p5p1/7p/8/5QPP/Pbq1rP2/3R2K1 w - - 3 30}, boardfontsize=11pt, pgfstyle=straightmove, arrow=to, linewidth=0.07em, color=blue!60!black, markmoves={d1-d7}, pgfstyle=straightmove, arrow=to, linewidth=0.07em, color=green!50!black, markmoves={d1-d8}]
}


\title{Can the review teach a move from his lessons instead of the engine's top line?\\[0.3em]
\large The lesson pick at the player's turning points}
\author{Chess Coach research lab}
\date{}

\begin{document}
\maketitle

\begin{abstract}
At each of the player's turning points the daily review teaches a better move and the line from
it. He asked for the move his study lessons teach, whenever the engine says it holds, named with
the lesson's rule, instead of the engine's top line. The app's lesson pick chooses among the
review's stored engine lines that hold for him: within \Window{} win\% of the engine's first line,
or keeping a won game won. Among those it takes the move a lesson teaches for a purpose his own
move lacked, reading each move's purpose with the board code. At \TrainTPs{} turning points up to
\TrainUntil{} and \HeldTPs{} after, a lesson taught the pick at \TrainCoverage\% and
\HeldCoverage\%. The pick gave up \TrainCost{} and \HeldCost{} win\% against the engine's first
line, and never left the holding verdicts. Players at his band choose it more often than the
engine's first line, by \TrainFind{} and \HeldFind{} points of Maia-2's probability. By the rule
fixed before the run, the pick is kept. It is the engine's first line at about \TrainIsTop\% of
turning points; where it is not, his band chooses it \TrainGainWhere{} and \HeldGainWhere{}
points more often. The lessons it names are few: \LessonsAll{} of the \Lessons, and two chapters
carry \TwoShare\% of the lesson picks, which shows where the next round should go.
\end{abstract}

\section{The question}

The daily review teaches, at each turning point, a better move and the line from it: the
engine's first line, or the move players at his band would likely play among those the engine
says hold. Since the app has study lessons (\Lessons{} of them, each a video and a written lesson
whose recommended moves were checked with Stockfish), the player asked for that move to come from
what the lessons teach whenever the engine says it holds: a queen trade, a rook to an open file,
a pawn break, a defence that meets the threat, named with the lesson's rule.

Two earlier studies bound what such a pick may claim. A rule of thumb explains a move but should
not choose it (rules-of-thumb), and a sound move he would find costs about what the engine says
it gives up (findable-move). So the lessons only choose among moves the engine already says hold.
At his turning points, how often does such a move exist, what does it cost against the engine's
first line, and do players at his band choose it more often than that first line?

\section{How it was measured}

\paragraph{The positions.} Every turning point in the app's archive with him to move and at least
one stored engine line: \Examined{} turning points. The \SkippedGames{} games of the two days read
while the change was designed are left out\SkippedLater{}. The candidates are the first moves of the stored
lines, up to \MaxCandidates{}, from the routine's deep search, each with its stored evaluation in
his view.
No new engine search was run.

\paragraph{What holds.} Each candidate is graded against the engine's first line, E, by the
review's own grading. It \emph{holds} for him when it is \emph{best}, within \Window{} win\% of
E (the review's window), or \emph{keeps} the result: E's position is won (at least \KeptWon{}
win\%, or a mate) and the candidate's still is.

\paragraph{His band.} The Maia-2 level of his rating at the game with the app's offset
(+\Offset), and a move's \emph{share} is Maia-2's probability for it at that band. N, the natural
move, is the holding candidate with the largest share.

\paragraph{The pick.} The board code reads what each move is for, from the move and its line:
it meets a threat, wins material, trades queens, takes an open file with a rook, breaks with a
pawn, improves the worst piece, and so on. The study service reads the same purposes from every
move the lessons recommend, and for a purpose offers the lesson chapters whose lesson is tagged
with one of its ideas, ranked by how well the chapter's own recommended moves match it. L, the
lesson pick, is the holding candidate his band chooses most often among those a lesson teaches
for a purpose his move lacked (\texttt{lesson}); else among those a lesson teaches
(\texttt{lesson-shared}); else the natural move; else E. His move's purposes are read with the
line the routine stored to punish it. The code is the app's at commit \texttt{\AppCommit}.

\paragraph{The rule, fixed before the run.} Train: the turning points up to \TrainUntil, the
first two thirds. Held out: the rest, up to \HeldoutUntil. On both, three bars. \emph{Coverage}:
a lesson teaches L (\texttt{lesson} or \texttt{lesson-shared}) at \CoverageBar\% of turning points
or more. \emph{Cost}: L's mean win\% loss against E is at most \CostBar, and no L falls outside the
holding verdicts. \emph{Findability}: L's share minus E's, paired by turning point, has its 95\%
interval above zero, from \Draws{} bootstrap draws over games (a game's turning points move
together). The pick is kept when all three pass on both.

\section{The verdict}

\RuleTable

All three bars pass on both splits (Table~\ref{tab:rule}), so the pick is kept. A lesson teaches
the pick at most turning points, and at most of them for a purpose his own move lacked. The pick
gives up \TrainCostMean{} and \HeldCostMean{} win\% on average. The largest losses, \TrainWorst{} and
\HeldWorst{} win\%, come from the \KeepsPicks{} picks that keep a won game won; every pick that
loses more than the window is one of them (\OverWindow{} in all). Players at his band choose the
pick a few points more often than the engine's first line, on both splits.

\section{What the pick looks like}

\PurposeTable

These numbers describe; none of them carries a rule.

\paragraph{Mostly the engine's move.} L is E at \TrainIsTop\% of train turning points and
\HeldIsTop\% of held-out ones: a lesson teaches the engine's own first line at \TrainTopTaught\%
and \HeldTopTaught\%. Where L is not E (\TrainDiffers{} and \HeldDiffers{} turning points), his
band chooses L \TrainGainWhere{} and \HeldGainWhere{} points more often than E. The average gain
in Table~\ref{tab:rule} is this gain spread over every turning point.

\paragraph{A little less natural than the natural move.} L is N at \TrainIsNatural\% of turning
points on both splits, and its share is \TrainVsNatural{} and \HeldVsNatural{} points below N's:
choosing the move a lesson teaches, rather than the one his band likes best, costs a little
findability.

\paragraph{Defence first.} The pick is most often taught for meeting a threat, \ThreatShare\% of
the lesson picks, then for a rook on an open file and for winning material
(Table~\ref{tab:purposes}).

\paragraph{Few chapters.} The \TaughtN{} lesson picks name \LessonsAll{} lessons and \ChaptersAll{}
chapters. One chapter carries \FirstShare\% of them: ``\FirstIdea'' from ``\FirstTitle'', named for
meeting a threat and for winning material alike, though \FirstNotQueen\% of its picks do not move
the queen. With ``\SecondIdea'' from ``\SecondTitle'' the two carry \TwoShare\%. The ranking reads
only the ideas a purpose shares with a chapter's tags, so one chapter stands for most defences.

\section{Four positions}

\begin{figure}[!ht]
\centering
\begin{subfigure}[t]{0.48\textwidth}
\centering
\EaBoard\\[0.4em]
{\footnotesize FEN \EaFen}
\caption{\EaSide{} to move. The pick (green) is \EaPick{} (\EaPickUci), chosen by \EaPShare\% of
his band; the engine's first line (blue) is \EaTop{} (\EaTopUci), at \EaTShare\%; the pick gives
up \EaLoss{} win\%. The board code: ``\EaSays''. The chapter: ``\EaIdea'' from ``\EaTitle''. He
played \EaPlayed{} (\EaPlayedUci).}
\label{fig:queens}
\end{subfigure}\hfill
\begin{subfigure}[t]{0.48\textwidth}
\centering
\EbBoard\\[0.4em]
{\footnotesize FEN \EbFen}
\caption{\EbSide{} to move. The pick is \EbPick{} (\EbPickUci), at \EbPShare\%; the engine's
first line is \EbTop{} (\EbTopUci), at \EbTShare\%; the pick gives up \EbLoss{} win\%. The board
code: ``\EbSays''. The chapter: ``\EbIdea'' from ``\EbTitle''. He played \EbPlayed{}
(\EbPlayedUci).}
\label{fig:piece}
\end{subfigure}\\[1em]
\begin{subfigure}[t]{0.48\textwidth}
\centering
\EcBoard\\[0.4em]
{\footnotesize FEN \EcFen}
\caption{\EcSide{} to move. The pick is \EcPick{} (\EcPickUci), at \EcPShare\%; the engine's
first line is \EcTop{} (\EcTopUci), at \EcTShare\%; the pick gives up \EcLoss{} win\%. The board
code: ``\EcSays''. The chapter: ``\EcIdea'' from ``\EcTitle''. He played \EcPlayed{}
(\EcPlayedUci).}
\label{fig:break}
\end{subfigure}\hfill
\begin{subfigure}[t]{0.48\textwidth}
\centering
\EdBoard\\[0.4em]
{\footnotesize FEN \EdFen}
\caption{\EdSide{} to move. The pick is \EdPick{} (\EdPickUci), at \EdPShare\%; the engine's
first line is \EdTop{} (\EdTopUci), at \EdTShare\%; the pick gives up \EdLoss{} win\%. The board
code: ``\EdSays''. The chapter named, ``\EdIdea'', moves the queen and takes a pawn; the pick
does neither. He played \EdPlayed{} (\EdPlayedUci).}
\label{fig:misfit}
\end{subfigure}
\caption{The pick in green, the engine's first line in blue, each with the share of his band
that chooses it and the chapter the study service ranks first for the pick's purpose. In the
last the pick is a defence and the lesson is on defence, but the chapter does not fit the move.}
\label{fig:examples}
\end{figure}

The four positions illustrate the summary; they test nothing. Each is the held-out turning point
where the pick, taught for a purpose his move lacked and not the engine's first line, gains most
share over that first line: among the queen trades, the worst piece improved or an outpost taken,
the pawn breaks, and the picks naming the most named chapter that do not move the queen.


\section{What the app can take}

\begin{itemize}[nosep]
\item The pick goes into the review's \texttt{teach}: among the moves that hold for him, the move
  a lesson teaches, with what it does on the board and the chapters that teach it. The daily video
  can then teach that move and its rule instead of the engine's verdict. Whether the video
  coaches better for it is a separate study (lesson-videos).
\item The chapter a pick names needs a fit check. Ranking chapters by the ideas their tags share
  with a purpose lets one chapter stand for most defences; a next round could rank them by how
  closely the chapter's own positions match this one: the piece that moves, what it captures, the
  threat it meets.
\item At \TrainNoLesson\% and \HeldNoLesson\% of turning points no holding move is one a lesson
  teaches. New purposes for the board code (an attack on the king,
  prophylaxis, a weak square) would let the lessons teach more of them.
\end{itemize}

\section{Caveats}

\begin{itemize}[nosep]
\item \textbf{Only the stored lines.} The candidates are the up to \MaxCandidates{} moves the
  review kept. The live \texttt{teach} also weighs his band's likeliest moves.
\item \textbf{The board code's purposes.} A move is taught for what the board code reads in it;
  a move can do more than its purposes say, or less.
\item \textbf{A lesson teaches a move} when a chapter's tags share an idea with the move's
  purpose. That says the lesson is on the right topic, not that the chapter's rule explains this
  position, as the last example shows.
\item \textbf{Share is not understanding.} Maia-2's share says how often players at his band
  choose a move, not whether he would follow the lesson behind it.
\item \textbf{The cost is the engine's.} In human hands a sound move he would find costs about
  what the engine says (findable-move); this study did not play the picks on.
\end{itemize}

\end{document}
